% registry: PI
\methodhead{PI}{Papuga and R\r{u}\v{z}i\v{c}ka (2008): integral criterion PI}\label{sheet:PI}

\block{Source} J.~Papuga, M.~R\r{u}\v{z}i\v{c}ka, \emph{Int.\ J.\ Fatigue} 30 (2008) 58--66
\cite{papuga2008}, Eq.~(16), Appendix~A (Eqs.~A.2, A.5); Table~2.

\block{Idea and scope} Stress-based integral fatigue-limit criterion, the companion of the
critical-plane criterion PCr of the same paper: the same kind of per-plane quantity, averaged over
all plane orientations instead of maximised. Proposed for the high-cycle fatigue limit of metals
under multiaxial loading.

\block{Inputs}
\begin{center}\small
\begin{tabular}{@{}llp{100mm}@{}}
\toprule
Symbol & Unit & Meaning\\\midrule
$\psi,\varphi$ & rad & polar and azimuth angles of the plane normal\\
$C_a$ & MPa & shear stress amplitude on the plane: radius of the minimum circumscribed circle of the shear path\\
$N_a$, $N_m$ & MPa & normal stress amplitude and mean on the plane\\
$f_{-1}=\fm$, $t_{-1}=\tm$ & MPa & fully reversed axial and torsion fatigue limits\\
$\kappa=f_{-1}/t_{-1}$ & -- & limit ratio\\
$a_I$, $b_I$ & --, MPa & weights, Eq.~\eqref{eq:pi-coef}\\
$\sigma_{H,m}$ & MPa & mean hydrostatic stress\\
$\sigma_a$, $\sigma_m$ & MPa & uniaxial amplitude and mean of the worked example\\
$\mathrm{DP}$ & -- & index, left-hand side of Eq.~\eqref{eq:pi-main} divided by $f_{-1}$\\
\bottomrule
\end{tabular}
\end{center}

\block{Formulation} Eq.~(16):
\begin{gather}
\sqrt{\frac{1}{4\pi}\int_0^{2\pi}\!\!\int_0^{\pi}
\left[a_I\,C_a^2+b_I\left(N_a+\kappa\,N_m\right)\right]\sin\psi\,d\psi\,d\varphi}\le f_{-1},
\label{eq:pi-main}\\
a_I=\tfrac52\,\kappa^2,\qquad b_I=f_{-1}\,(3-\kappa^2).
\label{eq:pi-coef}
\end{gather}
The orientation average of $N_m$ is $\sigma_{H,m}$, so the mean enters as $b_I\kappa\sigma_{H,m}$. In
uniaxial loading the averages are $\langle C_a^2\rangle=2\sigma_a^2/15$,
$\langle N_a\rangle=\sigma_a/3$, $\langle N_m\rangle=\sigma_m/3$, hence
\begin{equation}
(f_{-1}\,\mathrm{DP})^2=\frac{\kappa^2}{3}\,\sigma_a^2+\frac{f_{-1}(3-\kappa^2)}{3}\,(\sigma_a+\kappa\,\sigma_m).
\label{eq:pi-uni}
\end{equation}
Appendix~A (Eq.~A.2) integrates $N_a=t_{-1}\sin2\varphi\sin^2\psi$ in torsion as a \emph{signed}
quantity (integral 0). An amplitude is non-negative; with $\lvert N_a\rvert$ ($\langle\lvert N_a\rvert\rangle=4t_{-1}/(3\pi)$)
reversed torsion at $t_{-1}$ gives
\begin{equation}
\mathrm{DP}=\sqrt{1+(3-\kappa^2)\,\frac{4}{3\pi\kappa}} ,
\label{eq:pi-tor}
\end{equation}
not 1. The axial calibration (Eq.~A.5) holds.

\block{Constants and calibration} $a_I$, $b_I$ in closed form from $f_{-1}$ and $t_{-1}$
(Appendix~A). The repository uses the workbook \texttt{Sigm1}, \texttt{Taum1}; nothing is fitted
except the master-curve constants on the uniaxial $R=-1$ design tests.

\block{Life} None: a fatigue-limit criterion. The repository reads finite lives on the master curve
of the index (Section~\ref{sec:mastercurve}).

\block{Repository choices}
\begin{itemize}
\item Eq.~\eqref{eq:pi-main} as printed, with $N_a\ge0$. Reversed tension gives exactly 1; reversed
  torsion gives Eq.~\eqref{eq:pi-tor}. This property of the published criterion is kept visible,
  not re-calibrated.
\item Sphere average on a 48$\times$192 Gauss--Legendre/azimuth grid (weights summing to 1); the
  non-smooth $\lvert N_a\rvert$ integrates to about $4\cdot10^{-4}$ relative.
\item A negative radicand (compressive mean normal stress) is clipped to 0.
\item The index is the DP of the master curve. No \texttt{\_ext} variant.
\end{itemize}

\block{Worked example} Synthetic material: $f_{-1}=300$, $t_{-1}=187.5\MPa$ ($\kappa=1.6$). The index
is 1 in reversed tension at $f_{-1}$ (Eq.~A.5). Reversed torsion at $t_{-1}$: index 1.0567, equal to
Eq.~\eqref{eq:pi-tor}. Tension $\sigma_a=200$ with $\sigma_m=100\MPa$: index 0.7452, equal to
Eq.~\eqref{eq:pi-uni}.

\block{Verification} The paper prints no per-test values (its 119 tests are listed in Papuga's
thesis), only group statistics (Table~2). Checked: the parameters of Eq.~\eqref{eq:pi-coef}; the
Appendix~A sphere averages ($\langle C_a^2\rangle=\tfrac25t^2$ in torsion, $\tfrac{2}{15}f^2$ and
$\langle N_a\rangle=f/3$ in axial loading); the axial calibration and the torsion inconsistency;
the hydrostatic mean term; rotation invariance. The method audit reports that on the 43 tests of
Papadopoulos et al.\ (1997), Table~1, PI gives a mean of $-0.4\%$ and an SD of $5.0\%$, close to the
paper's Table~2; this comparison is not a test.
